\FloatBarrier
\section{Problem Definition}
\label{sec:problem}
The compiler must construct reusable RHS attribute programs and establish that
the completed rewrite is correct. Table~\ref{tab:notation} gives notation
shared by these requirements. Element types are fixed in this formulation.
\begin{table}[htbp]
\caption{Common notation for attribute synthesis and checking.}
\label{tab:notation}
\begin{tabularx}{\linewidth}{@{}lY@{}}\toprule
Symbol & Meaning \\\midrule
$s$ & Source information: shapes and operation attributes available to the attribute program.\\
$X$ & Input tensor values compatible with the shapes in $s$.\\
$a$ & RHS attributes to construct.\\
$\operatorname{LHS}(s,X)$ & Tensor computed by the LHS.\\
$\operatorname{RHS}(s,a,X)$ & Tensor computed by the RHS using attributes $a$.\\
$\operatorname{Guard}(s)$ & The user-written condition, evaluated successfully to \code{true}.\\\bottomrule
\end{tabularx}
\end{table}

\FloatBarrier
\subsection{What Must the Compiler Find?}
\label{sec:find}
Some RHS attributes are omitted or marked \code{auto}. Finding individual
attribute values at matching time would have the advantage of known ranks,
shapes, and source attributes. TEPL instead derives a formula before the
matched tensors are known. Its inputs are the user-defined operation semantics,
the rewrite rule, and any required conditions. Its output is a reusable program,
rather than an attribute value for one concrete tensor.

Formally, given the operation semantics, rewrite, and condition
$\operatorname{Guard}$, the compiler seeks an attribute program $f$ such that
\begin{equation}
 a=f(s).
 \label{eq:attribute-program}
\end{equation}
The program depends on source information, rather than tensor element values,
and must satisfy Section~\ref{sec:contract}'s contract for every legal source
instance satisfying the guard. For example, composing transposes with
permutations \code{inner} and \code{outer} gives the reusable formula
\code{result[i] = inner[outer[i]]}. Compilation derives the formula without
concrete permutations; matching supplies their values later. An explicit
constructor supplies $f$, while inference constructs it for omitted attributes.

\FloatBarrier
\subsection{What Makes the Result Correct?}
\label{sec:contract}
A completed replacement must have computable attributes, satisfy the RHS
operations' requirements, and produce the same result whenever the rule's
conditions hold. Both explicit constructors and inferred constructors require
checking for semantic rules. The required contract is
\begin{equation}
\begin{aligned}
\forall s.\quad &\operatorname{LegalLHS}(s)\land\operatorname{Guard}(s)\\
\Longrightarrow\quad
 &s\in\operatorname{dom}(f)\ \land\ \operatorname{LegalRHS}(s,f(s))\\
 &{}\land\ \forall X\text{ compatible with }s.\quad
 \operatorname{LHS}(s,X)=\operatorname{RHS}(s,f(s),X).
\end{aligned}
\label{eq:contract}
\end{equation}
In words, every legal LHS satisfying the conditions must yield a computable,
legal RHS with the same result for every compatible input value. Membership
in $\operatorname{dom}(f)$ requires, for example, that inverse permutations
and position lookups be evaluable. Equality includes ordered output shape
and element values under the selected arithmetic policy. Exact algebra alone
does not establish equality under floating-point reassociation.

The guard denotes successful evaluation to \code{true}; it is not an assumption
that a failing computation nevertheless returns a Boolean. The formal account
must respect ordered local evaluation, short-circuiting, indexing, checked
arithmetic, and evaluation failures. These details matter because metadata
programs can be partial. The compiler must establish the contract throughout
the legal guarded domain; otherwise it rejects the semantic rule as unproved.

\FloatBarrier
\subsection{When Can the Compiler Accept a Rewrite?}
The compiler accepts a semantic rewrite when the supported reasoning can show
that replacement attributes are computable, RHS operations are valid, and the
result is preserved under the stated conditions. The same standard applies
to explicit \code{derive} programs and inferred attributes. A rule rejected
as unproved may still be mathematically correct. It can be implemented
manually, but its correctness then has to be established separately.
The checks described below are intended to enforce
Equation~\eqref{eq:contract}; this draft does not promote that intended
contract to a proved soundness theorem for the implementation.
